\chapter{Raciocínio Lógico Matemático}
\label{cap:rlm}

\begin{edital}
Estruturas lógicas; lógica de argumentação (analogias, inferências,
deduções, conclusões); lógica sentencial (proposições, tabelas-verdade,
equivalências, diagramas); lógica de primeira ordem; raciocínio lógico
envolvendo problemas aritméticos, geométricos e matriciais.
\end{edital}
\peso{5 questões, peso 1.}

\begin{pegadinha}[Duas frentes: matemática E lógica formal --- não leia só uma]
Cuidado com uma leitura só. A \textbf{Dataprev 2024} foi quase toda
\textbf{conta} --- por isso a fama de ``RLM da FGV é matemática''. Mas o
\textbf{edital 2026} lista a lógica formal com peso próprio: estruturas
lógicas, lógica de argumentação, lógica sentencial e \textbf{lógica de
primeira ordem} (itens 1--4), e só no item 5 os ``problemas aritméticos,
geométricos e matriciais''. São só 5 questões, então não dá para prever o
mix --- \textbf{não pule a lógica formal apostando no padrão de 2024.} Cubra
as duas frentes.
\end{pegadinha}

\begin{center}
\begin{tabular}{@{}p{5cm}p{2cm}p{4.4cm}@{}}
\toprule
\textbf{Assunto} & \textbf{Prioridade} & \textbf{O que revisar} \\
\midrule
Porcentagem / aumentos sucessivos & altíssima & fator multiplicativo (1,30 $\times$ 1,10), desconto, lucro \\
Razão e proporção / regra de três & alta & direta e inversa, divisão proporcional \\
Análise combinatória & alta & arranjo, combinação, permutação, princípio multiplicativo \\
Médias e estatística & alta & média simples, ponderada, mediana, moda \\
Lógica: equivalências e argumentação & alta (edital 2026) & condicional, contrapositiva, De Morgan, validade \\
PA / PG & média & termo geral, soma \\
Geometria plana & média & área, perímetro, Pitágoras \\
Matrizes & média & operações, determinante \\
Diagramas lógicos / quantificadores & média & todo/algum/nenhum, primeira ordem \\
\bottomrule
\end{tabular}
\end{center}

\section{Ferramentas mínimas de matemática}

\begin{itemize}
\item \textbf{Aumentos sucessivos:} multiplique os fatores. +30\% depois
+10\% $\to$ 1,30 $\times$ 1,10 = 1,43 $\to$ aumento total \textbf{43\%} (não
40\%). Taxa média mensal: raiz, $\approx$ \textbf{19,6\%} (entre 19\% e
20\%).
\item \textbf{Divisão proporcional:} reparte na razão das partes (ex.:
prejuízo proporcional ao capital investido).
\item \textbf{Média ponderada:} $\Sigma(\text{nota} \times \text{peso}) /
\Sigma(\text{pesos})$. Colocar a menor nota no maior peso derruba a média ---
cuidado com ``necessariamente''.
\item \textbf{Combinação} $C(n,p) = n!/[p!(n-p)!]$ (ordem não importa)
$\times$ \textbf{arranjo} $A(n,p)$ (ordem importa). Grafo completo de $n$
vértices tem $C(n,2)$ arestas.
\item \textbf{Juros simples $\times$ compostos.} Simples: o juro incide
\textbf{sempre sobre o capital inicial}, $M = C(1 + i \cdot n)$. Compostos: o
juro incide \textbf{sobre o montante anterior} (juro sobre juro), $M = C(1 +
i)^n$. Em 2 períodos a diferença é exatamente o \emph{juro do juro do
primeiro período}. Ex.: R\$ 10.000 a 10\% a.a.\ por 2 anos $\to$ simples R\$
12.000, composto $10.000 \times 1{,}1^2 = $ R\$ 12.100; diferença R\$ 100.
\item \textbf{Progressão aritmética (PA):} termo geral $a_n = a_1 + (n-1)r$;
soma $S_n = (a_1 + a_n)\,n/2$. \textbf{PG:} $a_n = a_1 \cdot q^{\,n-1}$. O
distrator clássico usa $n$ no lugar de $n-1$ (com $a_1=7$ e $r=4$, o 10º
termo é $7 + 9{\cdot}4 = 43$, não $47$).
\item \textbf{Matrizes:} o produto $A \cdot B$ só existe se o número de
\textbf{colunas de $A$} for igual ao número de \textbf{linhas de $B$}; o
resultado tem as \textbf{linhas de $A$} e as \textbf{colunas de $B$}. Logo
$(2{\times}3) \cdot (3{\times}4) = (2{\times}4)$. O produto \textbf{não é
comutativo}.
\item \textbf{Conjuntos --- inclusão-exclusão:} $|A \cup B| = |A| + |B| - |A
\cap B|$. Quem está ``fora dos dois'' é o total menos a união. Ex.: 60
candidatos, 35 Java, 28 Python, 12 ambos $\to$ união $= 35+28-12 = 51$; fora
$= 60-51 = 9$. Esquecer de subtrair a interseção é o erro plantado.
\item \textbf{Medidas de posição:} \textbf{média} = soma/quantidade;
\textbf{mediana} = valor central do rol \textbf{ordenado} (com $n$ par, média
dos dois centrais); \textbf{moda} = o que mais se repete (pode não haver, ou
haver mais de uma). Ordenar antes de pegar a mediana é o passo que a banca
conta que você pule.
\item \textbf{Desvio padrão} mede a \textbf{dispersão em torno da média}, não
a posição da média. Duas séries com a \emph{mesma} média e desvios
diferentes: a de \textbf{menor} desvio é a mais \textbf{regular}
(concentrada). Desvio padrão não é amplitude --- não diz que os valores
variaram ``de 0 até o desvio''.
\end{itemize}

\section{Lógica formal (não subestime)}

\textbf{Conectivos e tabela-verdade:}

\begin{center}
\begin{tabular}{@{}p{4.2cm}p{1.6cm}p{4.2cm}@{}}
\toprule
\textbf{Conectivo} & \textbf{Símbolo} & \textbf{Falso quando} \\
\midrule
Conjunção ``e'' & $p \land q$ & ao menos um é falso \\
Disjunção ``ou'' & $p \lor q$ & ambos falsos \\
Condicional ``se...então'' & $p \to q$ & \textbf{V $\to$ F} (só nesse caso) \\
Bicondicional ``se e somente se'' & $p \leftrightarrow q$ & valores diferentes \\
Disjunção exclusiva ``ou...ou'' & $p \oplus q$ & valores iguais \\
\bottomrule
\end{tabular}
\end{center}

\textbf{Equivalências que a FGV cobra:}

\begin{itemize}
\item \textbf{Contrapositiva:} $p \to q \equiv \lnot q \to \lnot p$. ``Se é
fã então assiste'' $\equiv$ ``se \textbf{não} assiste então \textbf{não} é
fã''. (A recíproca $q \to p$ \textbf{não} é equivalente.)
\item \textbf{Condicional como disjunção:} $p \to q \equiv \lnot p \lor q$.
\item \textbf{Negação da condicional:} $\lnot(p \to q) \equiv p \land \lnot
q$.
\item \textbf{De Morgan:} $\lnot(p \land q) \equiv \lnot p \lor \lnot q$;
$\lnot(p \lor q) \equiv \lnot p \land \lnot q$.
\end{itemize}

\textbf{Argumentação (validade):} um argumento é \textbf{válido} se, sempre
que as premissas forem verdadeiras, a conclusão também for.

\begin{itemize}
\item \textbf{Modus ponens:} $p \to q, p \vdash q$ (válido).
\item \textbf{Modus tollens:} $p \to q, \lnot q \vdash \lnot p$ (válido).
\item \textbf{Falácias:} afirmar o consequente ($p \to q, q \vdash p$) e
negar o antecedente ($p \to q, \lnot p \vdash \lnot q$) são
\textbf{inválidas} --- a FGV usa como pegadinha.
\item \textbf{Silogismo hipotético:} $p \to q, q \to r \vdash p \to r$.
\end{itemize}

\textbf{Quantificadores / primeira ordem e diagramas:} ``Todo A é B'',
``Algum A é B'', ``Nenhum A é B''.

\begin{itemize}
\item Negações: negar ``todo A é B'' $\to$ ``\textbf{algum} A \textbf{não} é
B''; negar ``algum A é B'' $\to$ ``\textbf{nenhum} A é B''. A FGV adora a
negação errada do quantificador.
\item \textbf{Diagramas lógicos} (círculos de Euler/Venn): resolvem
``todo/algum/nenhum'' desenhando os conjuntos e testando o que
\textbf{necessariamente} decorre --- não o que ``pode ser''.
\end{itemize}

\begin{pegadinha}[Pegadinha recorrente de lógica]
Confundir \textbf{equivalência} com \textbf{implicação}, marcar a
\textbf{recíproca} como equivalente, ou negar quantificador trocando ``todo''
por ``nenhum'' (o certo é ``algum não'').
\end{pegadinha}

\begin{pegadinha}
\textbf{Nas de matemática:} cada alternativa errada é o resultado de um erro
de conta específico --- esqueceu de somar o peso, inverteu
numerador/denominador, contou o caso duplicado, somou as taxas em vez de
multiplicar, usou $n$ em vez de $n-1$ na PA, não subtraiu a interseção na
união de conjuntos, dividiu um grupo pelo outro em vez de pelo total na
probabilidade. Refaça com calma; o distrator ``quase certo'' vem de um passo
omitido.

\textbf{Nas de estatística:} o distrator confunde \textbf{dispersão} com
\textbf{posição} --- ``desvio padrão menor indica média menor'' é falso (a
média pode ser idêntica) --- ou trata o desvio como \textbf{amplitude}
(``variou entre 0 e 45 ms''). Na mediana, oferece o resultado de quem
\textbf{não ordenou} a série.

\textbf{Nas de lógica:} troca equivalência por implicação, marca a
recíproca como equivalente, nega o quantificador errado (``todo''$\to$
``nenhum'' em vez de ``algum não''), ou apresenta uma falácia (afirmar o
consequente) como argumento válido.
\end{pegadinha}

\begin{jacaiu}
Divisão proporcional (prejuízo); média ponderada com pesos por bimestre;
sistema com soma e soma dos quadrados; equivalência da condicional
(contrapositiva); grafo/estradas (combinação); aumentos sucessivos e taxa
média. Rode \texttt{./quiz.py rlm}.
\end{jacaiu}

\begin{comosair}
Chegue com as fórmulas na ponta: $C(n,2)=n(n-1)/2$; média ponderada =
$\Sigma(\text{nota}\times\text{peso})/\Sigma\text{pesos}$;
$(a+b)^2=a^2+2ab+b^2$; taxa acumulada = produto dos fatores; PA
$a_n=a_1+(n-1)r$; juros compostos $M=C(1+i)^n$ contra simples
$M=C(1+i{\cdot}n)$; $|A\cup B|=|A|+|B|-|A\cap B|$; produto de matrizes
$(m{\times}n)\cdot(n{\times}p)=(m{\times}p)$. Em porcentagem
sucessiva, \textbf{sempre multiplique} fatores; nunca some percentuais.
Estime a ordem de grandeza antes de marcar: isso elimina o distrator
``redondo'' plantado pela banca. Não gaste tempo demais numa questão de
conta pesada --- são só \textbf{5} no total (foram 6 na Dataprev 2024);
garanta as de porcentagem, média e combinatória, que são as mais
recorrentes.
\end{comosair}

\section{Alta probabilidade / pesquisa extra}

\begin{itemize}
\item \textbf{Probabilidade} básica: casos favoráveis / casos
\textbf{possíveis} (o total, não o outro grupo). Urna com 5 azuis e 3 verdes
$\to$ P(azul) $= 5/8$, não $5/3$ (nem probabilidade seria, passa de 1) nem
$5/3$ invertido em $3/5$.
\item \textbf{Diagramas lógicos} (todo/algum/nenhum) para argumentação.
\item Treine \textbf{velocidade}: são 5 questões valendo 1 ponto cada ---
não gaste tempo demais numa; os específicos valem 2,5$\times$.
\end{itemize}
